\documentclass[11pt]{article}

\usepackage[T1]{fontenc}
\usepackage{lmodern}
\usepackage[a4paper,margin=2.7cm]{geometry}
\usepackage{amsmath,amssymb,mathtools}
\usepackage{bm}
\usepackage{microtype}

\newcommand{\R}{\mathbb{R}}
\newcommand{\dt}{\,\mathrm{d}t}
\newcommand{\dx}{\,\mathrm{d}x}

\title{ADER--DG Formulation for the Exner Model}
\author{}
\date{}

\begin{document}

\maketitle

\section{The Exner model}

We consider the one-dimensional Exner model. Let \(h=h(x,t)\) denote the
water height, \(u=u(x,t)\) the depth-averaged water velocity,
\(q=q(x,t)=h(x,t)u(x,t)\) the water discharge, and \(b=b(x,t)\) the evolving
bottom topography. The free-surface elevation is
\[
    \eta=h+b,
\]
so that
\[
    h=\eta-b,
    \qquad
    u=\frac{q}{h}=\frac{q}{\eta-b}
\]
whenever \(h>0\).

The sediment-transport discharge is described by the classical Grass law
\[
    q_b=\xi A u\lvert u\rvert^{m_g-1},
    \qquad
    \xi=\frac{1}{1-\rho_0},
    \qquad
    0\leq A\leq 1,
    \qquad
    m_g=3,
\]
where \(\rho_0\in[0,1)\) is the porosity of the sediment layer. The empirical
coefficient \(A\) accounts for effects such as grain size and kinematic
viscosity; larger values of \(A\) correspond to stronger fluid--sediment
interaction.

The water mass equation and the Exner equation are, respectively,
\[
    \partial_t h+\partial_x q=0,
    \qquad
    \partial_t b+\partial_x q_b=0.
\]
Since \(\eta=h+b\), adding these two equations gives
\[
    \partial_t\eta+\partial_x(q+q_b)=0.
\]
Moreover, the momentum balance can be written in the pre-balanced form
\[
    \partial_t q
    +\partial_x\!\left(
        qu+\frac{g}{2}\bigl(\eta^2-2\eta b\bigr)
    \right)
    =-g\eta\,\partial_x b.
\]
Consequently, using the variables \((\eta,q,b)\), the Exner system takes the
compact form
\begin{equation}
    \partial_t\bm U(x,t)+\partial_x\bm F\bigl(\bm U(x,t)\bigr)
    =\bm S(x,t),
    \tag{6}
    \label{eq:exner-system}
\end{equation}
where
\[
    \bm U=
    \begin{pmatrix}
        \eta\\ q\\ b
    \end{pmatrix},
    \qquad
    \bm F(\bm U)=
    \begin{pmatrix}
        q+q_b\\[1mm]
        qu+\dfrac{g}{2}\bigl(\eta^2-2\eta b\bigr)\\[1mm]
        q_b
    \end{pmatrix},
    \qquad
    \bm S=
    \begin{pmatrix}
        0\\[1mm]
        -g\eta\,\partial_x b\\[1mm]
        0
    \end{pmatrix}.
\]
The third component \(q_b\) of the flux is required by the Exner equation
\(\partial_t b+\partial_x q_b=0\).

\section{ADER--DG formulation}

Unlike a Runge--Kutta DG discretization, the ADER--DG method uses the ADER
approach for the temporal evolution. The resulting method is one-step,
one-stage, and fully discrete.

\subsection{Mesh and approximation space}

Let the spatial domain be partitioned into \(N\) cells
\[
    I_j=\bigl[x_{j-\frac12},x_{j+\frac12}\bigr],
    \qquad j=1,\ldots,N,
\]
with cell centers and cell sizes
\[
    x_j=\frac12\left(x_{j-\frac12}+x_{j+\frac12}\right),
    \qquad
    \tau_j=x_{j+\frac12}-x_{j-\frac12},
    \qquad
    \tau=\max_{1\leq j\leq N}\tau_j.
\]
For the time interval \([t^n,t^{n+1}]\), define the local space--time control
volume
\[
    \Omega_j^n=I_j\times[t^n,t^{n+1}].
\]
The discontinuous finite-element space of degree at most \(k\) is
\[
    V_h^k
    =\left\{
        v:\ v\big|_{I_j}\in\mathbb P^k(I_j),
        \quad j=1,\ldots,N
      \right\}.
\]
Thus, the DG approximation satisfies
\(\bm U_h(\cdot,t)\in[V_h^k]^3\). Let
\(\phi_\ell=\phi_\ell(x)\in V_h^k\), \(\ell=0,\ldots,k\), be an arbitrary
time-independent test function.

\subsection{Space--time weak formulation}

Multiplying system~\eqref{eq:exner-system} by \(\phi_\ell\), integrating over
\(\Omega_j^n\), and integrating the temporal and spatial derivatives by parts
gives
\begin{align}
    &\int_{I_j}
        \bm U_{h,j}^{n+1}(x)\,\phi_\ell(x)\dx
    -
    \int_{I_j}
        \bm U_{h,j}^{n}(x)\,\phi_\ell(x)\dx     - \int_{t^n}^{t^{n+1}}\!\int_{I_j}
        \bm F\bigl(\bm U_{\tau,j}(x,t)\bigr)
        \partial_x\phi_\ell(x)\dx\dt
    \notag\\
    &\quad
    +
    \widehat{\bm F}_{j+\frac12}^{\,n}\,
        \phi_\ell\bigl(x_{j+\frac12}^{-}\bigr)
    -
    \widehat{\bm F}_{j-\frac12}^{\,n}\,
        \phi_\ell\bigl(x_{j-\frac12}^{+}\bigr) =
    \int_{t^n}^{t^{n+1}}\!\int_{I_j}
        \bm S_{\tau,j}(x,t)\,
        \phi_\ell(x)\dx\dt .
    \tag{7}
    \label{eq:ader-dg-weak-form}
\end{align}
This vector identity holds for \(j=1,\ldots,N\). Here,
\(\widehat{\bm F}_{j+\frac12}^{\,n}\) and
\(\widehat{\bm F}_{j-\frac12}^{\,n}\) denote the time-integrated numerical
flux contributions at the right and left interfaces of \(I_j\), respectively;
they may be constructed from a consistent Lax--Friedrichs numerical flux. Their
explicit computation is not considered at this stage.


\subsection{Initialization and local space--time prediction}

\paragraph{Step 1: Initialization.}

Let \(x_j\) denote the center of the cell \(I_j\). The initial DG
approximation in \(I_j\) is written as
\begin{equation}
    \mathbf{U}_{h,j}^{0}(x) = \mathbf{U}_{\tau,j}(x,0)
    =
    \sum_{k_x=0}^{k}
    \widetilde{\mathbf{U}}_j(k_x,0)
    (x-x_j)^{k_x},
    \qquad x\in I_j.
    \tag{9}
\end{equation}
The coefficients
\(\widetilde{\mathbf{U}}(k_x,0)\), \(k_x=0,\ldots,k\), are obtained
componentwise by the \(L^2\)-projection of the initial condition
\(\mathbf{U}_0(x)=\mathbf{U}(x,0)\) onto
\(\mathbb{P}^k(I_j)\). More precisely, they satisfy
\[
    \int_{I_j}
    \left(
        \mathbf{U}_{\tau}(x,0)-\mathbf{U}_0(x)
    \right)
    \phi_{\ell}(x)\,dx
    =
    \mathbf{0},
    \qquad \ell=0,\ldots,k.
\]

\paragraph{Step 2: Differential transformation procedure.}

Instead of employing the Cauchy--Kowalevski procedure used in
traditional ADER methods, we adopt the differential transformation
(DT) procedure. Suppose that the DG polynomial
\(\mathbf{U}_{h,j}^n(x)\) is known in the cell \(I_j\). The local
space--time approximation is expressed as the truncated Taylor
expansion
\begin{equation}
    \mathbf{U}_{\tau}(x,t)
    =
    \sum_{k_t=0}^{k}
    \sum_{k_x=0}^{k-k_t}
    \widetilde{\mathbf{U}}(k_x,k_t)
    (x-x_j)^{k_x}(t-t^n)^{k_t},
    \qquad (x,t)\in\Omega_j^n,
    \tag{10a}
\end{equation}
where
\begin{equation}
    \widetilde{\mathbf{U}}(k_x,k_t)
    =
    \frac{1}{k_x!\,k_t!}
    \left.
    \frac{\partial^{k_x+k_t}\mathbf{U}(x,t)}
         {\partial x^{k_x}\partial t^{k_t}}
    \right|_{(x,t)=(x_j,t^n)}
    \tag{10b}
\end{equation}
denotes the differential transform of \(\mathbf{U}\) at
\((x_j,t^n)\). The cell and time-step indices are omitted from the
transformed coefficients for simplicity.

The differential transformation rules needed in the present study
are given below. For arbitrary functions \(v(x,t)\), \(z(x,t)\), and
a constant \(c\), one has
\begin{align*}
    w(x,t)=c\,v(x,t)
    \quad &\Longrightarrow \quad
    \widetilde{w}(k_x,k_t)
    =
    c\,\widetilde{v}(k_x,k_t),                                      \\[0.3em]
    w(x,t)=\partial_x v(x,t)
    \quad &\Longrightarrow \quad
    \widetilde{w}(k_x,k_t)
    =
    (k_x+1)\widetilde{v}(k_x+1,k_t),                                \\[0.3em]
    w(x,t)=\partial_t v(x,t)
    \quad &\Longrightarrow \quad
    \widetilde{w}(k_x,k_t)
    =
    (k_t+1)\widetilde{v}(k_x,k_t+1),                                \\[0.3em]
    w(x,t)=v(x,t)z(x,t)
    \quad &\Longrightarrow \quad
    \widetilde{w}(k_x,k_t)
    =
    \sum_{r=0}^{k_x}\sum_{s=0}^{k_t}
    \widetilde{v}(r,s)
    \widetilde{z}(k_x-r,k_t-s).
\end{align*}
For \(w(x,t)=1/v(x,t)\), assuming
\(\widetilde{v}(0,0)\neq 0\), the transformed coefficients are
determined recursively by
\[
\left\{
\begin{aligned}
    \widetilde{w}(0,0)
    &=
    \frac{1}{\widetilde{v}(0,0)},                                   \\[0.4em]
    \widetilde{w}(k_x,k_t)
    &=
    -\frac{1}{\widetilde{v}(0,0)}
    \sum_{\substack{0\leq r\leq k_x,\;0\leq s\leq k_t\\r+s>0}}
    \widetilde{v}(r,s)
    \widetilde{w}(k_x-r,k_t-s),
    \qquad (k_x,k_t)\neq(0,0).
\end{aligned}
\right.
\]

Applying the DT procedure componentwise to
\[
    \partial_t\mathbf{U}
    +
    \partial_x\mathbf{F}(\mathbf{U})
    =
    \mathbf{S}
\]
gives
\begin{equation}
    \widetilde{\mathbf{U}}(k_x,k_t+1)
    =
    -\frac{k_x+1}{k_t+1}
    \widetilde{\mathbf{F}}(k_x+1,k_t)
    +
    \frac{1}{k_t+1}
    \widetilde{\mathbf{S}}(k_x,k_t).
    \tag{11}
\end{equation}
This recurrence is applied for
\[
    k_t=0,\ldots,k-1,
    \qquad
    k_x=0,\ldots,k-k_t-1.
\]

For the Exner system considered here, the transformed flux and source
are decomposed as
\begin{equation}
\left\{
\begin{aligned}
    \widetilde{\mathbf{F}}(k_x+1,k_t)
    &=
    \begin{pmatrix}
        \widetilde{Q}_1(k_x+1,k_t) \\[0.2em]
        \widetilde{G}_1(k_x+1,k_t)
        +\widetilde{G}_2(k_x+1,k_t)
        -\widetilde{G}_3(k_x+1,k_t) \\[0.2em]
        \widetilde{Q}_2(k_x+1,k_t)
    \end{pmatrix},
    \\[1em]
    \widetilde{\mathbf{S}}(k_x,k_t)
    &=
    \begin{pmatrix}
        0 \\[0.2em]
        -\widetilde{G}_4(k_x,k_t) \\[0.2em]
        0
    \end{pmatrix}.
\end{aligned}
\right.
\tag{12}
\end{equation}

We next express these intermediate quantities in terms of the
transformed coefficients of the main variables
\((\eta,q,b)\). Since
\[
    h=\eta-b,
    \qquad
    u=\frac{q}{h},
\]
we first define
\[
    \widetilde{h}(k_x,k_t)
    =
    \widetilde{\eta}(k_x,k_t)
    -
    \widetilde{b}(k_x,k_t).
\]
Let \(R=1/h\). Assuming \(\widetilde{h}(0,0)>0\), its transformed
coefficients satisfy
\[
\left\{
\begin{aligned}
    \widetilde{R}(0,0)
    &=
    \frac{1}{\widetilde{h}(0,0)},                                  \\[0.4em]
    \widetilde{R}(k_x,k_t)
    &=
    -\frac{1}{\widetilde{h}(0,0)}
    \sum_{\substack{0\leq r\leq k_x,\;0\leq s\leq k_t\\r+s>0}}
    \widetilde{h}(r,s)
    \widetilde{R}(k_x-r,k_t-s).
\end{aligned}
\right.
\]
Consequently,
\[
    \widetilde{u}(k_x,k_t)
    =
    \sum_{r=0}^{k_x}\sum_{s=0}^{k_t}
    \widetilde{q}(r,s)
    \widetilde{R}(k_x-r,k_t-s).
\]

Since \(m_g=3\), the Grass sediment transport law becomes
\[
    q_b=\xi A u|u|^2=\xi A u^3.
\]
Introducing
\[
    \widetilde{u^2}(k_x,k_t)
    =
    \sum_{r=0}^{k_x}\sum_{s=0}^{k_t}
    \widetilde{u}(r,s)
    \widetilde{u}(k_x-r,k_t-s),
\]
the transformed sediment discharge is
\begin{align*}
    \widetilde{Q}_2(k_x,k_t)
    &=
    \widetilde{q_b}(k_x,k_t)                                      \\
    &=
    \xi A
    \sum_{r=0}^{k_x}\sum_{s=0}^{k_t}
    \widetilde{u^2}(r,s)
    \widetilde{u}(k_x-r,k_t-s).
\end{align*}
Therefore, the first flux component is
\[
    \widetilde{Q}_1(k_x,k_t)
    =
    \widetilde{q}(k_x,k_t)
    +
    \widetilde{Q}_2(k_x,k_t).
\]

The transformed coefficients associated with the momentum flux are
\begin{align*}
    \widetilde{G}_1(k_x,k_t)
    &=
    \widetilde{(qu)}(k_x,k_t)                                     \\
    &=
    \sum_{r=0}^{k_x}\sum_{s=0}^{k_t}
    \widetilde{q}(r,s)
    \widetilde{u}(k_x-r,k_t-s),                                   \\[0.6em]
    \widetilde{G}_2(k_x,k_t)
    &=
    \frac{g}{2}\widetilde{(\eta^2)}(k_x,k_t)                       \\
    &=
    \frac{g}{2}
    \sum_{r=0}^{k_x}\sum_{s=0}^{k_t}
    \widetilde{\eta}(r,s)
    \widetilde{\eta}(k_x-r,k_t-s),                                \\[0.6em]
    \widetilde{G}_3(k_x,k_t)
    &=
    g\,\widetilde{(\eta b)}(k_x,k_t)                              \\
    &=
    g
    \sum_{r=0}^{k_x}\sum_{s=0}^{k_t}
    \widetilde{\eta}(r,s)
    \widetilde{b}(k_x-r,k_t-s).
\end{align*}

Finally, since the topography \(b=b(x,t)\) evolves in time, we have
\[
    \widetilde{b_x}(k_x,k_t)
    =
    (k_x+1)\widetilde{b}(k_x+1,k_t).
\]
It follows that
\begin{align*}
    \widetilde{G}_4(k_x,k_t)
    &=
    g\,\widetilde{(\eta b_x)}(k_x,k_t)                             \\
    &=
    g
    \sum_{r=0}^{k_x}\sum_{s=0}^{k_t}
    \widetilde{\eta}(r,s)
    \widetilde{b_x}(k_x-r,k_t-s)                                  \\
    &=
    g
    \sum_{r=0}^{k_x}\sum_{s=0}^{k_t}
    (k_x-r+1)
    \widetilde{\eta}(r,s)
    \widetilde{b}(k_x-r+1,k_t-s).
\end{align*}


Starting from the spatial expansion coefficients \(\widetilde{\mathbf{U}}(k_x,0)\), \(k_x=0,\ldots,k\), the recurrence relation~(11), together with the transformed expressions in~(12), is applied successively in increasing order of \(k_t\). This procedure determines all the space--time expansion coefficients

\[
    \widetilde{\mathbf{U}}(k_x,k_t),\qquad
    \widetilde{\mathbf{F}}(k_x,k_t),\qquad
    \widetilde{\mathbf{S}}(k_x,k_t),
    \qquad
    k_t=0,\ldots,k,\quad
    k_x=0,\ldots,k-k_t.
\]

\begin{align*}
    \mathbf{U}_{\tau,j}(x,t)
    &= \sum_{k_t=0}^{k} \sum_{k_x=0}^{k-k_t}
    \widetilde{\mathbf{U}}_j(k_x,k_t)
    (x-x_j)^{k_x}(t-t^n)^{k_t}, \\
    \mathbf{F}_{\tau,j}(x,t)
    &= \sum_{k_t=0}^{k} \sum_{k_x=0}^{k-k_t}
    \widetilde{\mathbf{F}}_j(k_x,k_t)
    (x-x_j)^{k_x}(t-t^n)^{k_t}, \\
    \mathbf{S}_{\tau,j}(x,t)
    &= \sum_{k_t=0}^{k} \sum_{k_x=0}^{k-k_t}
    \widetilde{\mathbf{S}}_j(k_x,k_t)
    (x-x_j)^{k_x}(t-t^n)^{k_t}.
\end{align*}
for \((x,t)\in\Omega_j^n\), \(j=1,\ldots,N\). In this manner, the local space--time evolution is obtained independently within each space--time control volume \(\Omega_j^n\).
 
\paragraph{Remark.}
The source coefficients \(\widetilde{\mathbf{S}}(k_x,k_t)\) are required  for
\(k_t=0,\ldots,k - 1,\quad k_x=0,\ldots,k-k_t-1\), which is precisely the range needed in recurrence~(11).


\subsection{Implementation steps of the ADER--DG method}

For clarity, the implementation of the proposed ADER--DG method over one
time step \(t^n\longrightarrow t^{n+1}\) can be summarized as follows.

\begin{enumerate}

    \item \textbf{Initialization.}
    
    At \(t=0\), compute the initial DG polynomial in every cell
    \(I_j\), \(j=1,\ldots,N\), by the componentwise \(L^2\)-projection
    of the initial data \(\mathbf{U}_0(x)=\mathbf{U}(x,0)\) onto
    \([\mathbb{P}^k(I_j)]^3\):
    \[
        \mathbf{U}_{h,j}^{0}(x)
        =
        \sum_{k_x=0}^{k}
        \widetilde{\mathbf{U}}_j(k_x,0)
        (x-x_j)^{k_x}.
    \]
    At a general time \(t^n\), the spatial coefficients
    \(\widetilde{\mathbf{U}}_j(k_x,0)\) are obtained from the known DG
    polynomial \(\mathbf{U}_{h,j}^{n}(x)\).

    \item \textbf{Local space--time prediction.}
    
    Starting from the coefficients
    \(\widetilde{\mathbf{U}}_j(k_x,0)\), apply the recurrence
    relation~\((11)\), together with the transformed flux and source
    expressions~\((12)\), successively in increasing order of \(k_t\).
    This gives
    \[
        \widetilde{\mathbf{U}}_j(k_x,k_t),\qquad
        \widetilde{\mathbf{F}}_j(k_x,k_t),\qquad
        \widetilde{\mathbf{S}}_j(k_x,k_t),
    \]
    for
    \[
        k_t=0,\ldots,k,
        \qquad
        k_x=0,\ldots,k-k_t.
    \]
    The corresponding local space--time predictors
    \(\mathbf{U}_{\tau,j}\), \(\mathbf{F}_{\tau,j}\), and
    \(\mathbf{S}_{\tau,j}\) are then reconstructed over
    \(\Omega_j^n=I_j\times[t^n,t^{n+1}]\). This predictor step is
    performed independently in every cell.

    \item \textbf{Evaluation of the interface fluxes.}
    
    At each interface \(x_{j+\frac12}\), evaluate the left and right
    predictor traces
    \[
        \mathbf{U}_{\tau,j}
        \bigl(x_{j+\frac12}^{-},t\bigr),
        \qquad
        \mathbf{U}_{\tau,j+1}
        \bigl(x_{j+\frac12}^{+},t\bigr),
    \]
    and use them to construct a consistent well-balanced numerical
    flux
    \[
        \mathbf{F}^{*}_{j+\frac12}(t)
        =
        \mathbf{F}^{*}\!\left(
            \mathbf{U}_{\tau,j}(x_{j+\frac12}^{-},t),
            \mathbf{U}_{\tau,j+1}(x_{j+\frac12}^{+},t)
        \right).
    \]
    Its time-integrated contribution is
    \[
        \widehat{\mathbf{F}}_{j+\frac12}^{\,n}
        =
        \int_{t^n}^{t^{n+1}}
        \mathbf{F}^{*}_{j+\frac12}(t)\,dt.
    \]
    The well-balanced interface flux must be coupled with a compatible
    discretization of the source term so that the targeted equilibrium
    state is preserved exactly.

    \item \textbf{Evaluation of the space--time integrals.}
    
    Evaluate all the volume, source, and interface integrals appearing
    in the weak formulation~\((7)\). In practice, these integrals are
    computed using sufficiently accurate Gaussian quadrature rules in
    both space and time. The quadrature order must be selected so that
    the polynomial terms involved in the ADER--DG discretization are
    integrated with the required accuracy.

    \item \textbf{One-stage DG update.}
    
    For each test function \(\phi_\ell\), \(\ell=0,\ldots,k\), compute
    the DG solution at \(t^{n+1}\) from the one-stage formula
    \[
    \begin{aligned}
        \int_{I_j}
        \mathbf{U}_{h,j}^{n+1}(x)\phi_\ell(x)\,dx
        &=
        \int_{I_j}
        \mathbf{U}_{h,j}^{n}(x)\phi_\ell(x)\,dx
        +
        \int_{t^n}^{t^{n+1}}\!\int_{I_j}
        \mathbf{F}_{\tau,j}(x,t)
        \partial_x\phi_\ell(x)\,dx\,dt
        \\
        &\quad
        -
        \widehat{\mathbf{F}}_{j+\frac12}^{\,n}
        \phi_\ell(x_{j+\frac12}^{-})
        +
        \widehat{\mathbf{F}}_{j-\frac12}^{\,n}
        \phi_\ell(x_{j-\frac12}^{+})
        \\
        &\quad
        +
        \int_{t^n}^{t^{n+1}}\!\int_{I_j}
        \mathbf{S}_{\tau,j}(x,t)
        \phi_\ell(x)\,dx\,dt .
    \end{aligned}
    \]
    The degrees of freedom of
    \(\mathbf{U}_{h,j}^{n+1}\) are obtained by solving the resulting
    local mass-matrix system.

    \item \textbf{Advance to the next time step.}
    
    Set \(n\leftarrow n+1\), determine the next time-step size according
    to the appropriate CFL restriction, and repeat Steps~2--5 until the
    final time is reached.

\end{enumerate}

\paragraph{Remark on nonlinear stabilization.}
A monotone interface flux, such as the local Lax--Friedrichs or Rusanov
flux, provides numerical dissipation at cell interfaces, but it does not,
by itself, prevent spurious oscillations in the higher-order DG polynomial
near shocks or strong discontinuities. Therefore, for shock-dominated
solutions, a TVB, WENO, or another suitable nonlinear limiter should be
used. The limiter is applied to the updated spatial DG polynomial
\(\mathbf{U}_{h,j}^{n+1}\), preferably only in cells detected as troubled,
and not directly to the local space--time predictor
\(\mathbf{U}_{\tau,j}(x,t)\). The limiting procedure must preserve the
cell average and should be designed consistently with the well-balanced
property. As an alternative
for very strong shocks, a troubled-cell ADER--DG method with a robust
first-order finite-volume numerical fluxes where needed may be employed.


\subsection{DG Approximation and Local Space--Time Predictor}

Throughout the following presentation, we distinguish between the spatial
DG approximation and the local space--time predictor. At the discrete time
$t^n$, the approximate DG solution is denoted by $\mathbf{U}_h^n(x)$, whose
restriction
\[
    \mathbf{U}_{h,j}^n(x)
    =
    \left.\mathbf{U}_h^n(x)\right|_{I_j}
\]
is a polynomial of degree at most $k$ in the spatial variable.

Starting from $\mathbf{U}_{h,j}^n$, the differential transformation
procedure constructs, independently in each cell $I_j$, a local
space--time predictor $\mathbf{U}_{\tau,j}(x,t)$ over
$I_j\times[t^n,t^{n+1}]$. This predictor is initialized by the known DG
polynomial:
\begin{equation}
    \mathbf{U}_{\tau,j}(x,t^n)
    =
    \mathbf{U}_{h,j}^n(x),
    \qquad x\in I_j.
    \label{eq:predictor-initialization}
\end{equation}
Indeed, the local predictor is written as
\[
    \mathbf{U}_{\tau,j}(x,t)
    =
    \sum_{k_t=0}^{k}
    \sum_{k_x=0}^{k-k_t}
    \widetilde{\mathbf{U}}_j(k_x,k_t)
    (x-x_j)^{k_x}(t-t^n)^{k_t}.
\]
Evaluating this expression at $t=t^n$ eliminates all terms with
$k_t\geq 1$ and gives
\[
    \mathbf{U}_{\tau,j}(x,t^n)
    =
    \sum_{k_x=0}^{k}
    \widetilde{\mathbf{U}}_j(k_x,0)(x-x_j)^{k_x}
    =
    \mathbf{U}_{h,j}^n(x).
\]
Therefore, the coefficients
$\widetilde{\mathbf{U}}_j(k_x,0)$ are precisely the coefficients of the
known DG polynomial $\mathbf{U}_{h,j}^n$ expressed in the monomial basis
centered at $x_j$.

At the initial time $t^0=0$, these coefficients are obtained from the
$L^2$-projection of the prescribed initial condition $\mathbf{U}_0$ onto
$\mathbb{P}^k(I_j)$. At every subsequent time $t^n$, $n\geq 1$, they are
obtained directly from the previously computed DG polynomial
$\mathbf{U}_{h,j}^n$; thus, no new projection of the initial condition is
performed.

The predictor $\mathbf{U}_{\tau,j}$ is an auxiliary polynomial used to
evaluate the space--time volume integrals, source terms, and time-dependent
interface fluxes appearing in the ADER--DG formulation. Since it is
constructed locally and independently in each cell, it does not yet
include the numerical-flux corrections associated with interactions
between neighboring cells. Consequently, although
\[
    \mathbf{U}_{\tau,j}(x,t^n)
    =
    \mathbf{U}_{h,j}^n(x),
\]
one generally has
\[
    \mathbf{U}_{\tau,j}(x,t^{n+1})
    \neq
    \mathbf{U}_{h,j}^{n+1}(x).
\]
The updated DG solution $\mathbf{U}_h^{n+1}$ is obtained only after applying
the globally coupled DG correction step, in which the local predictors are
used to compute the volume contributions and numerical fluxes. The
ADER--DG evolution can therefore be summarized as
\[
    \mathbf{U}_h^n
    \;\longrightarrow\;
    \mathbf{U}_{\tau}
    \;\longrightarrow\;
    \mathbf{U}_h^{n+1},
\]
where the first arrow represents the element-local space--time prediction
and the second represents the globally coupled DG update.

\end{document}